% ==============================================================================
% BAB I: PENDAHULUAN
% ==============================================================================

\chapter{PENDAHULUAN}

\section{Latar Belakang}
Evolusi arsitektur \textit{web service} modern telah bergeser secara masif dari peladen fisik (\textit{bare-metal}) dan mesin virtual (\textit{Virtual Machine} / VM) menuju ekosistem kontainerisasi berbasis Docker \citep{khan2026decomposing, gurramkonda2026docker, nurzaman2022analisis}. Kontainerisasi menawarkan portabilitas tinggi, efisiensi konsumsi sumber daya komputasi, serta kecepatan \textit{deployment} aplikasi berskala besar. Namun, dalam lingkungan produksi yang sesungguhnya, kluster kontainer \textit{backend} jarang beroperasi dalam kondisi homogen yang ideal. Pembatasan kuota komputasi melalui fitur \textit{Control Groups} (cgroups) pada kernel Linux, fragmentasi perangkat keras, serta alokasi dinamis pada arsitektur \textit{microservices} menciptakan asimetri kapasitas komputasi (baik alokasi vCPU maupun kuota RAM) antar kontainer \textit{backend} \citep{shetty2025cpu, faridah2026performance}. Dalam praktik industri komputasi awan, variasi kapasitas ini umumnya mengadopsi pola tingkatan instansi berjenjang (\textit{instance tiering} atau \textit{t-shirt sizing}: \textit{Small}, \textit{Medium}, \textit{Large}) dengan rasio penggandaan kapasitas eksponensial seperti 1:2:4, yang mencerminkan heterogenitas terencana antar node layanan.

Penerapan penyeimbang beban (\textit{load balancer}) seperti NGINX dan HAProxy telah menjadi standar industri dalam mendistribusikan lalu lintas HTTP menuju jajaran peladen aplikasi \citep{buke2026load, sakib2025dynamic, arini2025perbandingan}. Kendati demikian, kedua perangkat lunak penyeimbang beban tersebut dibangun di atas filosofi arsitektur internal yang berbeda secara fundamental. NGINX mengadopsi model arsitektur \textit{multi-process} berbasis \textit{master-worker}, di mana sekumpulan proses pekerja (\textit{worker}) mengeksekusi \textit{non-blocking event-loop} secara mandiri dalam ruang memori terpisah. Sebaliknya, HAProxy dirancang murni sebagai proksi lapis aplikasi berkinerja tinggi yang beroperasi dengan model proses tunggal berorientasi peristiwa (\textit{event-driven single process}) dan memanfaatkan struktur data pohon biner elastis (\textit{elastic binary tree} / \texttt{ebtree}) berkecepatan $O(1)$ untuk manajemen antrean koneksi serta penjadwalan waktu respons deterministik \citep{tarreau2024haproxy}. Perbedaan arsitektural konkurensi ini diprediksi menghasilkan respons stabilitas yang sangat berlainan ketika penyeimbang beban berhadapan dengan kluster kontainer \textit{backend} heterogen, terlebih saat menerima karakteristik beban kerja yang bertolak belakang, yaitu beban intensif komputasi (\textit{CPU-bound}) versus beban intensif masukan/keluaran (\textit{I/O-bound}) \citep{baeturohman2024analisis}.

Mayoritas penelitian terdahulu yang membandingkan performa NGINX dan HAProxy pada lingkungan kontainer Docker hanya mengevaluasi kondisi \textit{backend} yang homogen serta terpaku pada algoritma statis tanpa bobot seperti \textit{Round Robin} \citep{nurzaman2022analisis, ramadhani2026analysis}. Di sisi lain, penelitian yang menguji algoritma berbobot (\textit{Weighted Round Robin} dan \textit{Weighted Least Connections}) pada lingkungan heterogen umumnya diimplementasikan pada Virtual Private Server (VPS) atau mesin virtual terisolasi tanpa mempertimbangkan karakteristik isolasi kernel cgroups Docker \citep{rizqi2022analisis}. Ketidaksesuaian antara logika algoritma distribusi beban dan batasan kuota cgroups berpotensi memicu fenomena \textit{CPU throttling} akibat pembatasan Completely Fair Scheduler (CFS) kernel Linux \citep{shetty2025cpu}, penumpukan antrean koneksi (\textit{backlog}), dan lonjakan ekstrem pada latensi ekor (\textit{tail latency} persentil p95 dan p99) \citep{wodislawski2026ponytail, li2025tailbench}. Oleh karena itu, diperlukan penelitian komparatif yang menguji secara empiris dan terukur kinerja data-plane NGINX versus HAProxy pada lingkungan Docker dengan kapasitas heterogen.

\section{Rumusan Masalah}
Berdasarkan latar belakang masalah yang telah diuraikan, rumusan masalah dalam penelitian ini adalah sebagai berikut:
\begin{enumerate}
    \item Bagaimana perbandingan kinerja data-plane NGINX dan HAProxy menggunakan algoritma \textit{Round Robin}, \textit{Least Connections}, \textit{Weighted Round Robin}, dan \textit{Weighted Least Connections} pada kluster kontainer Docker dengan kapasitas \textit{backend} heterogen?
    \item Bagaimana pengaruh karakteristik beban kerja intensif prosesor (\textit{CPU-bound}) dan beban kerja intensif masukan/keluaran (\textit{I/O-bound}) terhadap \textit{throughput}, latensi rata-rata, dan \textit{tail latency} (p95 dan p99) pada masing-masing mesin penyeimbang beban?
    \item Sejauh mana presisi distribusi beban aktual yang dihasilkan oleh algoritma penyeimbang beban terhadap proporsi kapasitas kuota cgroups kontainer \textit{backend}?
\end{enumerate}

\section{Batasan Masalah}
Agar penelitian ini terarah dan mencapai sasaran yang diharapkan, batasan masalah ditetapkan sebagai berikut:
\begin{enumerate}
    \item Perangkat lunak penyeimbang beban yang dievaluasi dibatasi pada NGINX \textit{Open Source} (versi stabil) dan HAProxy \textit{Community Edition}.
    \item Algoritma penyeimbang beban yang diuji meliputi empat varian: \textit{Round Robin} (RR), \textit{Least Connections} (LC), \textit{Weighted Round Robin} (WRR), dan \textit{Weighted Least Connections} (WLC).
    \item Lingkungan uji dibangun dengan kontainer Docker yang dikelola melalui Docker Compose pada sistem operasi Linux (kernel 6.x).
    \item Heterogenitas kapasitas kontainer \textit{backend} disimulasikan menggunakan fitur cgroups dengan perbandingan rasio kapasitas vCPU dan RAM sebesar 1:2:4 (masing-masing 0.5 vCPU/256 MB, 1.0 vCPU/512 MB, dan 2.0 vCPU/1024 MB) yang merepresentasikan tingkatan ukuran instansi standar (\textit{cloud compute tiers}: \textit{Small}, \textit{Medium}, \textit{Large}).
    \item Karakteristik pengujian beban dibagi menjadi dua kategori: \textit{endpoint} CPU-bound (komputasi intensif) dan \textit{endpoint} I/O-bound (penundaan I/O asinkron).
    \item Alat uji beban (\textit{load testing tool}) menggunakan Grafana k6 dengan metrik utama: \textit{throughput} (req/s), \textit{mean latency}, \textit{p95/p99 tail latency}, dan \textit{error rate}.
\end{enumerate}

\section{Tujuan Penelitian}
Tujuan yang hendak dicapai dalam pelaksanaan penelitian ini adalah:
\begin{enumerate}
    \item Menganalisis dan membandingkan performa NGINX dan HAProxy dalam menangani trafik HTTP pada lingkungan kontainer Docker berkapasitas heterogen.
    \item Mengevaluasi efektivitas algoritma penjadwalan berbobot (WRR dan WLC) dibandingkan algoritma tanpa bobot (RR dan LC) terhadap pencegahan lonjakan \textit{tail latency} dan \textit{CPU throttling}.
    \item Mengukur presisi distribusi trafik aktual untuk membuktikan kesesuaian antara alokasi bobot algoritma penyeimbang beban dengan kapasitas nyata kontainer \textit{backend}.
\end{enumerate}

\section{Manfaat Penelitian}
Hasil dari penelitian ini diharapkan dapat memberikan manfaat sebagai berikut:
\begin{enumerate}
    \item \textbf{Manfaat Teoretis}: Memberikan kontribusi ilmiah dan literatur empiris mengenai interaksi kausal antara mesin proksi lapis data-plane (arsitektur \textit{multi-process} NGINX vs \textit{event-driven} HAProxy) dengan mekanisme isolasi cgroups kernel Linux pada kluster kontainer heterogen.
    \item \textbf{Manfaat Praktis}: Menjadi panduan teknis bagi insinyur infrastruktur, praktisi \textit{DevOps}, dan arsitek sistem dalam memilih perangkat lunak penyeimbang beban serta menentukan algoritma distribusi yang optimal pada ekosistem komputasi awan berbasis kontainer.
\end{enumerate}

\section{Sistematika Penulisan}
Sistematika penulisan laporan proposal skripsi ini disusun dalam tiga bab utama sebagai berikut:
\begin{itemize}
    \item \textbf{BAB I PENDAHULUAN}: Memaparkan latar belakang masalah, perumusan masalah, batasan masalah, tujuan penelitian, manfaat penelitian, serta sistematika penulisan.
    \item \textbf{BAB II TINJAUAN PUSTAKA}: Menguraikan kajian literatur penelitian terdahulu yang relevan, dasar teori pendukung, serta kedudukan dan kebaruan penelitian.
    \item \textbf{BAB III METODOLOGI PENELITIAN}: Memaparkan alur kerangka penelitian, perancangan arsitektur topologi sistem, spesifikasi sumber daya heterogen, matriks skenario pengujian, instrumen evaluasi data, serta jadwal pelaksanaan penelitian.
\end{itemize}
